\chapter{Sequences and Limits}

A small program prints one number per line. The rule is only one: the next term is
the square root of two plus the previous term, starting from one. What appears on
the screen: 1.000000, then 1.732051, 1.931852, 1.982890, 1.995718, 1.998929, and
so on. On the twelfth line the screen shows 1.999999738545. That number appears to
stop moving. But appearing and being true are two different things: the program
does not know whether the thousandth line will still shift, and it also does not
know which number this run is heading toward.

The second example is more mischievous. The sequence $(1+1/n)^n$ for $n=10$ has
value 2.5937424601, for $n=100$ has value 2.7048138294, for $n=1000$ has value
2.7169239322, and for $n=10^7$ has value 2.7182816941. The first six digits appear
settled at 2.71828. But the true value at $n=10^7$ is 2.7182816925, so the last
seven digits in that machine output are rounding error, not a sign of closeness.
A table can deceive in two directions at once: it makes a divergent sequence look
calm, and it makes a convergent sequence look off course. Chapter 22 will explain
why machine arithmetic behaves that way; this chapter prepares the tools that let
us not depend on tables.

Chapter 3 already supplied the raw material: the field axioms, the order axioms,
and the supremum property. But the supremum only works once the set of terms has
been collected; it does not tell us whether the sequence stops moving. What we
need is a definition that can be checked on paper, and then theorems that turn
that check into estimated work: the $\varepsilon$-$N$ definition, the algebraic
properties of limits, the squeeze theorem, monotone sequences, subsequences, the
Cauchy criterion, and the limit superior.

\begin{untukanak}
Have you ever walked toward a door under a strange rule: each step covers only
half of the distance that remains? The first step takes you halfway. The second
cuts half of what is left, so you are already three quarters of the way. The third
adds another eighth. Every step moves you forward, yet there is always a remainder
you have not reached. If that rule is kept up, you never actually touch the door.

Even so, there are questions that can be answered with certainty. How many steps
are needed so that the distance to the door is less than a hand's width? The answer
can be computed before you set out, and it does not change when the room becomes
twice as long: you only need a few extra steps. So there are two different things
here. ``Never arrives'' and ``as close as one likes can be approached'' are not two
contradictory statements. In fact both are true at once.

Now the question that makes one curious. What if the rule is changed slightly, say
each step covers a third of the remaining distance, or half plus a little? Can you
still guess where that journey ends, and can its speed still be computed? In this
chapter we will know why such questions have definite answers, and why we never
obtain the answers merely by watching numbers on a screen.
\end{untukanak}

\section{The definition of a limit and how to use it}
\tingkat{basic}

\kamus{sequence}{barisan}{a run of numbers given consecutive numbers, like the house
numbers along one side of a street that never ends}

A limit answers one question: how close the later terms gather, and from when on.
In everyday language, a sequence tends to $L$ when, after a few initial terms are
discarded, all the remaining terms are as close as one likes to $L$.

\begin{definisi}[Convergence]
A sequence $(x_n)$ of real numbers \emph{converges} to $L \in \mathbb{R}$ if for
every $\varepsilon>0$ there is $N \in \mathbb{N}$ such that for every $n>N$ we have
$|x_n-L|<\varepsilon$. We write $x_n \to L$, and $L$ is called the
\emph{limit}\index{limit of a sequence} of the sequence. A sequence that does not
converge to any real number is called \emph{divergent}\index{divergent sequence}.
\end{definisi}

The order of the quantifiers in that definition decides everything. The value
$\varepsilon$ is chosen first, then $N$ follows and may depend on $\varepsilon$.
Reversing that order produces a far stronger and usually false statement: there are
sequences that satisfy ``for every $n$ there is $N$'' vacuously, because $N=n$ can
always be used. A healthy way to read it is as a game between two parties. The
challenger names an $\varepsilon$, however small. The responder must answer with an
$N$ that works for all later terms. The sequence converges when the responder
always has an answer, for every challenge.

Figure~\ref{fig:pita-epsilon} captures that game for $x_n=1+1/n$. Two bands centered
at $L=1$ are drawn with different widths, and each has its own threshold: $N=2$ for
$\varepsilon=0.5$, $N=10$ for $\varepsilon=0.1$. Nothing about the sequence changes;
only the challenger's demand changes.

\begin{figure}[htbp]
\centering
\begin{tikzpicture}[x=0.72cm,y=3.0cm,>=Stealth,font=\footnotesize]
  \fill[layerteal!14] (0,0.5) rectangle (14,1.5);
  \fill[layerorange!20] (0,0.9) rectangle (14,1.1);
  \draw[->,ink] (-0.1,0.35) -- (14.6,0.35) node[right] {$n$};
  \draw[->,ink] (0,0.35) -- (0,2.15) node[above] {$x_n$};
  \draw[dashed,layergrey,thick] (0,1) -- (14.1,1);
  \node[layergrey,anchor=south west] at (0.08,1.02) {$L=1$};
  \foreach \n/\y in {1/2.0000,2/1.5000,3/1.3333,4/1.2500,5/1.2000,6/1.1667,%
      7/1.1429,8/1.1250,9/1.1111,10/1.1000,11/1.0909,12/1.0833,13/1.0769,14/1.0714}
    \fill[accent] (\n,\y) circle (2.2pt);
  \draw[dashed,layerteal,thick] (2,0.35) -- (2,1.52);
  \draw[dashed,layerorange,thick] (10,0.35) -- (10,1.12);
  \node[anchor=north] at (2,0.30) {$N_1=2$};
  \node[anchor=north] at (10,0.30) {$N_2=10$};
  \node[anchor=north west,align=left] at (4.6,2.28)
    {outer band: $\varepsilon=0.5$\\ inner band: $\varepsilon=0.1$};
\end{tikzpicture}
\caption{The sequence $x_n=1+1/n$ entering the $\varepsilon$ band around $L=1$. With
$\varepsilon=0.5$ every term after $n=2$ is already inside the band; with
$\varepsilon=0.1$ the threshold shifts to $n=10$. A narrower band demands a more
distant threshold.}
\label{fig:pita-epsilon}
\end{figure}

\begin{bacakeras}
Two voices take turns. The first asks, the second answers. Read it aloud slowly,
pause at every period, and do not add words in between.

``Longer and longer?'' \ ``Closer and closer.'' \ ``How close?'' \ ``As close as
requested.'' \ ``What if one thousandth is requested?'' \ ``Wait a thousand steps,
then measure.'' \ ``What if one hundredth of a millionth?'' \ ``Wait one hundred
million steps. Patience costs nothing.'' \ ``So it will arrive?'' \ ``No. It never
arrives.'' \ ``Then what is the use of waiting?'' \ ``Because the distance shrinks,
and the one who asks is free to choose how small.'' \ ``And if the requests stop?''
\ ``Then there is nothing to prove.''

The line repeated three times, more slowly each time: close is not arrived; close
is a distance we are allowed to set ourselves.
\end{bacakeras}

For $x_n=1/n$ the answer is immediate: take $N$ a natural number with
$N>1/\varepsilon$, which exists by the Archimedean property in Chapter 3. For $n>N$
we have $|1/n-0|=1/n<1/N<\varepsilon$. If the challenger asks for
$\varepsilon=10^{-3}$, the answer is $N=1000$; at $n=1000$ the error is exactly
$10^{-3}$ and does not yet meet the strict condition, whereas at $n=1001$ the error
is $0.000999$. This one-term difference is not pedantry: the definition uses $n>N$,
not $n \ge N$, so the $N$ that is chosen must genuinely pass the threshold.

\begin{tangan}
Materials: duct tape, a tape measure, chalk, and one sheet of notes. Mark a number
line on the floor two meters long: zero at the left end, two at the right end, and
the number one exactly in the middle. Stick down two horizontal strips as wide as a
foot, centered at the number one: a wide band from one half to one and a half
meters, a narrow band from nine tenths to one point one meters.

Now walk to the position one plus one over n, for n equal to one, two, three, up to
fourteen. Each time, note whether your heel lands in the wide band, in the narrow
band, or outside both, then walk back so the marks are not stepped on.

What should be visible: the first two steps are still outside the narrow band, the
tenth step stops exactly at its edge, and after that no heel comes out again. A
narrower band demands more steps, while the number one is never touched by a heel.
\end{tangan}

The use of a definition is felt most strongly when we want to prove that something
does \emph{not} converge. The negation of the definition reads: there is
$\varepsilon>0$ such that for every $N$ there is $n>N$ with $|x_n-L| \ge
\varepsilon$. Note that $L$ is still a free name there; to refute convergence we
must do it for every $L$.

\begin{contoh}
The sequence $x_n=(-1)^n$ is divergent.
\end{contoh}

\begin{proof}
Suppose $x_n \to L$ for some $L$. Take $\varepsilon=1$; there is $N$ such that
$|x_n-L|<1$ for all $n>N$. Choose an even $n>N$; then $x_n=1$ and $|1-L|<1$. The
next term $n+1$ also exceeds $N$ and is odd, so $x_{n+1}=-1$ and $|-1-L|<1$. The
triangle inequality gives
\[
2 = |1-(-1)| \le |1-L| + |L-(-1)| < 1+1 = 2,
\]
an impossibility. So no $L$ satisfies the definition.
\end{proof}

The proof just given shows a pattern that will appear often: to refute, we choose an
$\varepsilon$ smaller than the distance between two groups of terms. The even and
odd groups are a fixed distance $2$ apart, so no single $N$ can separate them when
$\varepsilon<1$.

The first property worth recording is uniqueness. In everyday words: a sequence
cannot stop in two different places at once.

\begin{teorema}[Uniqueness of the limit]
If $x_n \to L$ and $x_n \to M$, then $L=M$.
\end{teorema}

\begin{proof}
Suppose $L \neq M$ and let $\varepsilon=|L-M|/2>0$. There is $N_1$ with
$|x_n-L|<\varepsilon$ for $n>N_1$, and $N_2$ with $|x_n-M|<\varepsilon$ for
$n>N_2$. For $n>\max(N_1,N_2)$ both inequalities hold at once, so
\[
|L-M| \le |L-x_n| + |x_n-M| < 2\varepsilon = |L-M| .
\]
The inequality $|L-M|<|L-M|$ is impossible, so the supposition is false.
\end{proof}

The phrase ``for $n$ sufficiently large'' means: there is $N$ such that the stated
property holds for all $n>N$. With that vocabulary, the definition of a limit in
brief: $x_n$ is as close as one likes to $L$ when $n$ is sufficiently large.

The second property: a convergent sequence cannot run off to infinity.

\begin{teorema}
Every convergent sequence is bounded.
\end{teorema}

\begin{proof}
Suppose $x_n \to L$. Take $\varepsilon=1$; there is $N$ with $|x_n-L|<1$ for
$n>N$, so $|x_n| \le |x_n-L|+|L| < |L|+1$. Only the first term up to term $N$ is
still uncontrolled, and there are finitely many of those. Hence
\[
M=\max\{|x_1|,|x_2|,\dots,|x_N|,|L|+1\}
\]
is a bound that works for every $n$.
\end{proof}

The converse does not hold: $(-1)^n$ is bounded but divergent. Boundedness is a
necessary condition, not a sufficient one; from a table that looks calm we may only
conclude boundedness.

\section{Algebraic properties of limits}
\tingkat{intermediate}

In everyday language: a limit is polite toward addition, multiplication, and
division, as long as the divisor does not tend to zero.

Before proving those properties we need one small control: a sequence tending to a
nonzero number eventually stays away from zero.

\begin{lemma}
If $y_n \to y$ with $y \neq 0$, then there is $N$ such that $|y_n|>|y|/2$ for all
$n>N$.
\end{lemma}

\begin{proof}
Take $\varepsilon=|y|/2>0$. There is $N$ with $|y_n-y|<|y|/2$ for $n>N$. Then
$|y_n| \ge |y|-|y_n-y| > |y|-|y|/2 = |y|/2$.
\end{proof}

The consequence for computation: limits can be computed term by term, and the only
exception is division by zero.

\begin{teorema}[Algebraic properties of limits]
If $x_n \to x$ and $y_n \to y$, then
$x_n+y_n \to x+y$, \ $x_n y_n \to xy$, and $c x_n \to cx$ for every
$c \in \mathbb{R}$. If $y \neq 0$ and $y_n \neq 0$ for all $n$, then
$x_n/y_n \to x/y$.
\end{teorema}

\begin{proof}
For the sum, take $\varepsilon>0$. Choose $N_1$ with $|x_n-x|<\varepsilon/2$ for
$n>N_1$ and $N_2$ with $|y_n-y|<\varepsilon/2$ for $n>N_2$. For
$n>\max(N_1,N_2)$,
$|(x_n+y_n)-(x+y)| \le |x_n-x|+|y_n-y| < \varepsilon$. The scalar multiple case is
similar and shorter: for $c=0$ the sequence is constant; for $c \neq 0$ use
$\varepsilon/|c|$.

For the product, write the difference as two parts,
\[
x_n y_n - xy = x_n(y_n-y) + (x_n-x)y .
\]
Because $(x_n)$ converges it is bounded: there is $M>0$ with $|x_n| \le M$ for all
$n$. Choose $N$ so that for $n>N$ we have $|y_n-y|<\varepsilon/(2M)$ and
$|x_n-x|<\varepsilon/(2(|y|+1))$. Then
\[
|x_n y_n - xy| \le M|y_n-y| + |x_n-x|\,|y| < \frac{\varepsilon}{2}
+ \frac{\varepsilon}{2} = \varepsilon .
\]
The denominator $|y|+1$ is used so that it is never zero when $y=0$; it stays a
positive constant, so it does not change the conclusion.

For the quotient it suffices to show $1/y_n \to 1/y$, then use the product result.
By the lemma, for $n>N_1$ we have $|y_n|>|y|/2$, so
\[
\left|\frac{1}{y_n}-\frac{1}{y}\right| = \frac{|y-y_n|}{|y_n y|}
\le \frac{2}{|y|^2}\,|y_n-y| .
\]
The right side is smaller than $\varepsilon$ as long as $|y_n-y|<\varepsilon|y|^2/2$,
and that can be guaranteed by choosing $N_2$.
\end{proof}

Note the technique that has just appeared: adding and subtracting the same term to
split a difference into two small parts. Almost every convergence proof involving a
product uses it, and the threshold taken is always the largest one.

The next property concerns order, not operations. In everyday language: a limit
cannot jump over a fence that all terms obey.

\begin{teorema}[Order property]
If $x_n \to x$ and $x_n \ge 0$ for all $n$, then $x \ge 0$. More generally, if
$x_n \le y_n$ for all sufficiently large $n$, $x_n \to x$, and $y_n \to y$, then
$x \le y$.
\end{teorema}

\begin{proof}
Suppose $x<0$ and take $\varepsilon=-x>0$. There is $N$ with $|x_n-x|<\varepsilon$
for $n>N$, so in particular $x_n < x+\varepsilon = 0$. That contradicts
$x_n \ge 0$. For the second statement, apply the first part to $y_n-x_n \ge 0$,
whose limit is $y-x$ by the previous theorem.
\end{proof}

Note that the legitimate conclusion is $x \ge 0$, not $x>0$: the sequence $1/n$ is
entirely positive but its limit is zero. The order property also must not be used
backwards. From $x_n<y_n$ we only obtain $x \le y$; a strict gap can vanish in the
limit.

With the algebraic properties, limits of rational sequences can be computed as
usual. For
\[
x_n=\frac{3n^2+2n}{n^2-5},
\]
divide numerator and denominator by $n^2$:
$x_n=\dfrac{3+2/n}{1-5/n^2}$. Because $1/n \to 0$ and $1/n^2 \to 0$, the numerator
tends to $3$ and the denominator to $1 \neq 0$, so $x_n \to 3$. The numerical values
do creep there: $x_{10}=3.368421$, $x_{100}=3.021511$. Another example,
\[
y_n=\frac{2n^2-n}{3n^2+7} \longrightarrow \frac{2}{3},
\]
with $y_{10}=0.618893$ and $y_{100}=0.663179$; both are still far from $2/3$ at
$n=100$ because the error is about $1/n$.

The algebraic properties speak only about limits that already exist; applying them
to a sequence not yet known to converge is an expensive mistake.

\section{The squeeze theorem and comparison sequences}
\tingkat{intermediate}

Sometimes the limit of a sequence cannot be computed directly, but it can be trapped
from above and from below by two sequences with the same limit. In everyday
language: when two guards walk close together with a child between them, the child
must stop in the same place as they do.

\begin{teorema}[Squeeze theorem]
\index{squeeze theorem}
Suppose $a_n \le x_n \le b_n$ for all sufficiently large $n$. If
$a_n \to L$ and $b_n \to L$, then $x_n \to L$.
\end{teorema}

\begin{proof}
Take $\varepsilon>0$. Choose $N_1$ so that $|a_n-L|<\varepsilon$ for $n>N_1$, that
is $a_n>L-\varepsilon$, and $N_2$ so that $b_n<L+\varepsilon$ for $n>N_2$. Also take
$N_3$ so that the squeezing holds for $n>N_3$. For $n>\max(N_1,N_2,N_3)$,
\[
L-\varepsilon < a_n \le x_n \le b_n < L+\varepsilon,
\]
so $|x_n-L|<\varepsilon$.
\end{proof}

Before using it we need one comparison sequence that will keep appearing: powers that
shrink to zero.

\begin{lemma}
If $0<q<1$, then $q^n \to 0$.
\end{lemma}

\begin{proof}
Write $1/q=1+h$ with $h>0$. Bernoulli's inequality, proved by induction in Chapter
1, gives $(1+h)^n \ge 1+nh$. Since both sides are positive,
\[
0 < q^n = \frac{1}{(1+h)^n} \le \frac{1}{1+nh} .
\]
Given $\varepsilon>0$, the Archimedean property gives $n$ with $1+nh>1/\varepsilon$,
that is $1/(1+nh)<\varepsilon$. By the squeeze theorem, $q^n \to 0$.
\end{proof}

Now an example that cannot be settled by dividing by the highest power. For
$x_n=(2^n+3^n)^{1/n}$, the largest term inside the brackets is $3^n$; factor that
power out:
\[
x_n = 3\left(1+\left(\tfrac{2}{3}\right)^{n}\right)^{1/n}.
\]
Because $(2/3)^n \to 0$ and $1 \le 1+(2/3)^n \le 2$, we obtain
$1 \le x_n/3 \le 2^{1/n}$. It remains to show $2^{1/n} \to 1$: write
$2^{1/n}=1+a_n$ with $a_n \ge 0$, then
$2=(1+a_n)^n \ge 1+na_n$, so $0 \le a_n \le 1/n$. The squeeze theorem closes both
ends, and $x_n \to 3$. Python numbers confirm the direction:
$x_1=5$, $x_2=3.605551$, $x_5=3.075152$, $x_{10}=3.005162$, and $x_{100}=3$.

\begin{untukanak}
Two school guards walk side by side along a corridor, and a pupil walks between them.
The guard on the left steps toward the exit, the guard on the right does too, and the
two get closer until their shoulders almost touch. The pupil squeezed in the middle
does not need to know the way; he only needs to stay between the two guards. Wherever
those guards stop, there the pupil stops as well.

Now the question that makes one curious. What if the left guard stops earlier than
the right guard: can we still guess where the pupil stops? And what if their distance
is never zero, only smaller from day to day? Does the pupil still have a definite
place to stop, or does he keep walking with no end?

It turns out the answer does not depend on the faces of the guards, only on two
things: where the left guard is heading, and where the right guard is heading. If
both head for the same place, the pupil in the middle has no other choice. That way
of thinking, ``squeezed between two parties'', is what mathematicians use for
sequences whose formula is too complicated to compute directly.
\end{untukanak}

\section{Monotone sequences and the monotone convergence theorem}
\tingkat{intermediate}

In everyday language: a sequence that moves in only one direction, and never crosses
a certain fence, must stop somewhere. This is where completeness works most cheaply,
because the limit need not be guessed in advance.

\kamus{monotone sequence}{barisan monoton}{a staircase whose steps only go up, never
down}

The everyday language is simple: one direction only, never turning back.

\begin{definisi}[Monotone sequence]
A sequence $(x_n)$ is called \emph{increasing}\index{monotone sequence} if
$x_n \le x_{n+1}$ for all $n$, and \emph{decreasing} if $x_n \ge x_{n+1}$ for all
$n$. A sequence is called monotone if it is increasing or decreasing.
\end{definisi}

A rough picture: what moves in one direction without crossing a fence must stop, and
it stops exactly at the nearest fence it never crosses.

\begin{teorema}[Monotone convergence]
\index{monotone convergence theorem}
An increasing sequence bounded above converges, and its limit is
$\sup\{x_n : n \in \mathbb{N}\}$.
\end{teorema}

\begin{proof}
Suppose $(x_n)$ is increasing and bounded above. Its set of terms is nonempty and
bounded above, so by the supremum property in Chapter 3 it has a supremum; call it
$s=\sup\{x_n\}$. Take $\varepsilon>0$. Because $s-\varepsilon$ is not an upper bound,
there is $N$ with $x_N>s-\varepsilon$. For $n>N$, monotonicity gives
$x_n \ge x_N>s-\varepsilon$, while $x_n \le s$ by the definition of supremum. So
$|x_n-s|<\varepsilon$.
\end{proof}

The immediate consequence: for such a sequence the limit can be read off from its set
of terms, without computing term by term.

\begin{akibat}
A decreasing sequence bounded below converges, and its limit is the infimum of its
set of terms.
\end{akibat}

\begin{proof}
Apply the theorem to $(-x_n)$, which is increasing and bounded above, then use the
scalar multiple property of limits and the identity
$\sup\{-x_n\}=-\inf\{x_n\}$.
\end{proof}

Note that this theorem does not say the limit equals the bound we use. The bound $3$
in the following example is only a proof device; the limit is in fact $e$.

The first application is the recursive sequence that opens this chapter. Suppose
$x_1=1$ and $x_{n+1}=\sqrt{2+x_n}$.

First, all terms lie between $1$ and $2$. For $n=1$ this is clear. If
$1 \le x_n \le 2$, then $3 \le 2+x_n \le 4$, so
$1 \le \sqrt{3} \le x_{n+1} \le 2$. By induction, $1 \le x_n \le 2$ for all $n$.

Second, the sequence is increasing. Because $x_n \le 2$,
\[
x_{n+1} \ge x_n \iff 2+x_n \ge x_n^2 \iff (2-x_n)(1+x_n) \ge 0,
\]
and the last inequality is true since $x_n \le 2$ and $x_n \ge 1$.

By the monotone convergence theorem, $(x_n)$ converges; call its limit $L$. Because
$x_n \le 2$ for all $n$, the order property gives $L \le 2$. Now square the relation:
$x_{n+1}^2=2+x_n$. The left side tends to $L^2$ by the product property, the right
side to $2+L$, so $L^2=2+L$, or $(L-2)(L+1)=0$. Since $L \ge 1$, the only
possibility is $L=2$.

Note that we never use continuity of the square root; the algebraic properties of
limits and the order property are enough, so the machinery of Chapter 11 is not
needed here.

The second application is the sequence $y_n=(1+1/n)^n$. The binomial expansion gives
\[
y_n = \sum_{k=0}^{n} \binom{n}{k}\frac{1}{n^k}
= \sum_{k=0}^{n} \frac{1}{k!}\prod_{j=1}^{k-1}\left(1-\frac{j}{n}\right),
\]
with the empty product equal to $1$ for $k=0$ and $k=1$. Every factor $1-j/n$
increases as $n$ grows, so every term with index $k$ increases; the sum also gains
one new positive term. Hence $(y_n)$ is increasing. Every factor is at most $1$, so
\[
y_n \le \sum_{k=0}^{n}\frac{1}{k!} \le 1+1+\frac12+\frac14+\cdots = 3,
\]
because $1/k! \le 1/2^{k-1}$ for $k \ge 1$, which is checked by induction. The
sequence is increasing and bounded above by $3$, so it converges. We give its limit
the name $e$.

Figure~\ref{fig:monoton-e} displays that sequence for $n=1$ up to $n=30$. The points
rise without ever turning back, the dashed line at $3$ is the fence used in the proof,
and the dotted line at $2.71828$ is where it actually stops. The gap between fence and
destination stays there throughout the figure; the fence works because it is never
crossed, not because it is reached.

\begin{figure}[htbp]
\centering
\begin{tikzpicture}[x=0.32cm,y=3.0cm,>=Stealth,font=\footnotesize]
  \draw[->,ink] (-0.6,1.85) -- (31.2,1.85) node[right] {$n$};
  \draw[->,ink] (0,1.85) -- (0,3.12) node[above] {$y_n$};
  \draw[dashed,layergrey,thick] (0,3.0) -- (30.6,3.0);
  \node[layergrey,anchor=south] at (20,3.02) {proof bound $B=3$};
  \draw[densely dotted,layerindigo,very thick] (0,2.71828) -- (30.4,2.71828);
  \node[layerindigo,anchor=west] at (31.3,2.71828) {limit $e\approx2.71828$};
  \foreach \n/\y in {1/2.0000,2/2.2500,3/2.3704,4/2.4414,5/2.4883,6/2.5216,%
      7/2.5465,8/2.5658,9/2.5812,10/2.5937,11/2.6042,12/2.6130,13/2.6206,%
      14/2.6272,15/2.6329,16/2.6379,17/2.6424,18/2.6464,19/2.6500,20/2.6533,%
      21/2.6563,22/2.6590,23/2.6615,24/2.6637,25/2.6658,26/2.6678,27/2.6696,%
      28/2.6713,29/2.6728,30/2.6743}
    \fill[accent] (\n,\y) circle (1.7pt);
  \draw[->,layerteal,thick] (9.5,2.42) -- (15.5,2.60);
  \node[layerteal,anchor=north] at (12,2.40) {increasing, never turning back};
\end{tikzpicture}
\caption{The sequence $y_n=(1+1/n)^n$ for $n=1$ up to $n=30$. The points rise and
never turn back; the fence $B=3$ is never touched; the value approached is
$e\approx2.71828$, not the fence.}
\label{fig:monoton-e}
\end{figure}

\section{Subsequences, Bolzano--Weierstrass, and the Cauchy criterion}
\tingkat{advanced}

In everyday language: from a bounded sequence one can always pick a run of children
that stops somewhere, even when the parent run is chaotic.

\kamus{subsequence}{subbarisan}{sorting out some letters from a pile without changing
their order}

A subsequence is formed by choosing only some terms: the indices used grow larger,
but any number of terms may be skipped.

\begin{definisi}[Subsequence]
Let $n_1<n_2<n_3<\cdots$ be increasing natural numbers. The sequence $(x_{n_k})_{k
\in \mathbb{N}}$ is called a \emph{subsequence}\index{subsequence} of $(x_n)$.
\end{definisi}

Because $n_k \ge k$ for every $k$, which is checked by induction, a subsequence
cannot be ``slower'' than its parent sequence in terms of index.

In everyday language: choosing some terms does not change the stopping place, it only
slows down or speeds up its arrival.

\begin{proposisi}
If $x_n \to L$, then every subsequence $(x_{n_k})$ also converges to $L$.
\end{proposisi}

\begin{proof}
Take $\varepsilon>0$ and $N$ with $|x_n-L|<\varepsilon$ for $n>N$. For $k>N$ we have
$n_k \ge k>N$, so $|x_{n_k}-L|<\varepsilon$.
\end{proof}

The consequence is simple but most often used to refute: two different stopping places
are already enough.

\begin{akibat}
If $(x_n)$ has two subsequences with different limits, then $(x_n)$ is divergent.
\end{akibat}

\begin{proof}
If $(x_n)$ converged, all its subsequences would converge to the same limit, and
uniqueness of the limit forbids two different values.
\end{proof}

This corollary explains the example $(-1)^n$ without an argument by impossibility: the
even subsequence tends to $1$, the odd subsequence to $-1$. The way it works also
explains why divergence is more often provable than convergence.

To prove Bolzano--Weierstrass we need the following lemma. Its content: every
sequence, however chaotic, contains a monotone subsequence.

\begin{lemma}
Every real sequence has a monotone subsequence.
\end{lemma}

\begin{proof}
Call an index $n$ a \emph{peak} if $x_n \ge x_m$ for all $m>n$, that is, if no later
term is larger. There are two possibilities.

If there are infinitely many peaks, list them as $n_1<n_2<\cdots$. Because $n_1$ is a
peak, $x_{n_1} \ge x_{n_2}$; because $n_2$ is a peak, $x_{n_2} \ge x_{n_3}$; and so
on. The result is a decreasing subsequence.

If there are only finitely many peaks, take $n_1$ beyond the last peak. Because $n_1$
is not a peak, there is $n_2>n_1$ with $x_{n_2}>x_{n_1}$. Because $n_2$ also lies
beyond the last peak, it is not a peak either, so there is $n_3>n_2$ with
$x_{n_3}>x_{n_2}$. The process continues without end and produces an increasing
subsequence.
\end{proof}

The content of this theorem: from a trapped sequence one can always pick some terms
that genuinely tend to a single point.

\begin{teorema}[Bolzano--Weierstrass]
\index{bolzano-weierstrass}
Every bounded real sequence has a convergent subsequence.
\end{teorema}

\begin{proof}
By the lemma, $(x_n)$ has a monotone subsequence. A subsequence of a bounded sequence
is also bounded. A bounded monotone sequence converges by the monotone convergence
theorem. Hence that subsequence converges.
\end{proof}

An alternative proof uses repeated bisection of intervals and the nested interval
property; both proofs lean on completeness, which we borrowed above through the
supremum.

\kamus{Cauchy sequence}{barisan Cauchy}{passengers crowding ever tighter in a train
corridor, without knowing where the exit door is}

Now consider a different problem. The definition of a limit requires us to name $L$
first, whereas in many computations it is $L$ that we are seeking. All we have is a
list of numbers. We therefore need a measure of closeness between the terms
themselves, without naming a destination.

\begin{definisi}[Cauchy sequence]
A sequence $(x_n)$ is called \emph{Cauchy}\index{cauchy sequence} if for every
$\varepsilon>0$ there is $N$ such that for all $m,n>N$ we have
$|x_n-x_m|<\varepsilon$.
\end{definisi}

Note the difference in the order of ideas. In the definition of a limit all terms are
compared with one fixed number $L$; in the Cauchy definition terms are compared with
other terms, with no outside number. That is why the Cauchy definition makes sense in
any space that has a distance, whereas the limit definition demands a candidate
destination.

\begin{teorema}[Cauchy criterion]
\index{cauchy sequence}
A real sequence converges if and only if it is Cauchy.
\end{teorema}

\begin{proof}
Suppose $x_n \to L$ and take $\varepsilon>0$. There is $N$ with
$|x_n-L|<\varepsilon/2$ for $n>N$. For $m,n>N$,
\[
|x_n-x_m| \le |x_n-L| + |L-x_m| < \varepsilon .
\]
So convergence gives Cauchy.

Conversely, suppose $(x_n)$ is Cauchy. Take $\varepsilon=1$; there is $N_0$ such that
$|x_n-x_m|<1$ for $m,n>N_0$. Choosing $m=N_0+1$ gives
$|x_n|<|x_{N_0+1}|+1$ for $n>N_0$, so $(x_n)$ is bounded. By Bolzano--Weierstrass
there is a subsequence $x_{n_k} \to L$. Take $\varepsilon>0$; choose $N_1$ with
$|x_n-x_m|<\varepsilon/2$ for $m,n>N_1$, and $K$ with $|x_{n_k}-L|<\varepsilon/2$ for
$k>K$. Choose $k>K$ that also satisfies $n_k>N_1$; then for $n>N_1$,
\[
|x_n-L| \le |x_n-x_{n_k}| + |x_{n_k}-L| < \varepsilon .
\]
So $x_n \to L$.
\end{proof}

The second direction uses Bolzano--Weierstrass, which uses the monotone convergence
theorem, which uses the supremum. That chain cannot be cut: inside $\mathbb{Q}$ this
theorem fails. The sequence of decimal truncations
$1;\ 1.4;\ 1.41;\ 1.414;\ 1.4142;\ \dots$ tending to the square root of two is Cauchy,
since the difference of two terms is at most $10^{-k}$ when both pass term $k$, yet no
rational number is its limit. If that sequence converged to $L$, then $a_k^2 \to L^2$;
on the other hand $a_k^2 \to 2$, so $L^2=2$. Chapter 3 already showed that no rational
has square $2$. So ``Cauchy'' and ``convergent'' are two different properties; what
makes them agree is completeness, and that is the content of the supremum axiom.

\begin{tantangan}
The sequence $h_n=1+\frac12+\cdots+\frac1n$ grows very slowly: $h_{10}=2.928968$,
$h_{100}=5.187378$, $h_{1000}=7.485471$, and $h_{10^6}=14.392727$. The difference of
two consecutive terms, namely $1/(n+1)$, tends to zero, so this sequence passes the
shallowest test, the one that most often deceives: its terms get closer together. The
question is whether $(h_n)$ converges. Prove your answer, and explain why
$h_{10^6}<15$ is not a proof that the sequence is bounded.

Hint 1: do not inspect terms one by one. Sum in blocks whose lengths double, that is
one term, two terms, four terms, eight terms, and so on.

Hint 2: the $k$-th block contains $2^{k-1}$ terms, and the smallest term in it is at
least $1/2^k$; so every block contributes at least $1/2$. Compare the result with
$h_{2^m}$ and see what happens as $m$ grows.

If you are stuck, return to the Monotone Convergence Theorem: there an increasing
sequence is proved convergent precisely because a fence exists, and here it is
precisely the fence that is missing.
\end{tantangan}

\section{Limit superior and limit inferior}
\tingkat{advanced}

In everyday language: for a sequence that oscillates, we can still name the highest and
the lowest place it oscillates to, even though the sequence itself has no limit.

\kamus{limit superior}{limit superior}{the highest peak still visited by the sequence
as often as one likes after passing any point}

Throughout this section $(x_n)$ is bounded. For every $n$ form the tail set
$\{x_k : k \ge n\}$ and write
\[
s_n=\sup\{x_k : k \ge n\}, \qquad t_n=\inf\{x_k : k \ge n\}.
\]
Because the tails shrink, $(s_n)$ is decreasing and $(t_n)$ is increasing; both are
bounded, so both converge by the corollary to the monotone convergence theorem.

In everyday language: what we track is not where the sequence is, but how high and how
low it still visits after any point.

\begin{definisi}[Limit superior and inferior]
For a bounded sequence $(x_n)$,
\[
\limsup_{n\to\infty} x_n = \lim_{n\to\infty} s_n
\qquad \text{and} \qquad
\liminf_{n\to\infty} x_n = \lim_{n\to\infty} t_n .
\]
The first number is called the \emph{limit superior}\index{limit superior}, the second
the \emph{limit inferior}\index{limit inferior}.
\end{definisi}

Because $s_n \ge t_n$ for every $n$ and limits preserve order, we always have
$\liminf x_n \le \limsup x_n$. The two can differ greatly, and precisely that
difference is the measure of the chaos of the sequence.

\begin{contoh}
For $x_n=(-1)^n+1/n$, the tail contains the even terms $1+1/k$ and the odd terms
$-1+1/k$. The supremum of the $n$-th tail comes from the smallest even term still
available, so $s_1=1.5$, $s_2=1.5$, $s_3=1.25$, $s_4=1.25$, $s_5=1.166667$, and so on
decreasing toward $1$. The infimum of the tail behaves differently: the odd terms
$-1+1/k$ are always present in every tail and decrease toward $-1$ without ever
reaching it, so $t_n=-1$ for all $n$, even though not a single term equals $-1$. Hence
$\limsup x_n=1$ and $\liminf x_n=-1$, while the sequence is divergent.
\end{contoh}

Figure~\ref{fig:divergen-ayun} plots the $n$-th term against $n$ for the same sequence.
The run of even points clings near $1$, the run of odd points clings near $-1$, and the
distance between the two groups is fixed at $2$. No single horizontal line can be
approached by all the points, however far we walk.

\begin{figure}[htbp]
\centering
\begin{tikzpicture}[x=0.62cm,y=1.9cm,>=Stealth,font=\footnotesize]
  \draw[->,ink] (-0.3,0) -- (13.4,0) node[right] {$n$};
  \draw[->,ink] (0,-1.3) -- (0,2.0) node[above] {$x_n$};
  \draw[densely dotted,layerindigo,very thick] (0,1) -- (13.0,1);
  \draw[densely dotted,layerplum,very thick] (0,-1) -- (13.0,-1);
  \node[layerindigo,anchor=south west] at (0.1,1.03) {$\limsup x_n=1$};
  \node[layerplum,anchor=north west] at (0.1,-1.03) {$\liminf x_n=-1$};
  \foreach \n/\y in {1/0.0000,2/1.5000,3/-0.6667,4/1.2500,5/-0.8000,6/1.1667,%
      7/-0.8571,8/1.1250,9/-0.8889,10/1.1000,11/-0.9091,12/1.0833}
    \fill[accent] (\n,\y) circle (2.2pt);
  \draw[<->,layerbrown,thick] (12.5,-1) -- (12.5,1) node[midway,right] {distance $2$};
\end{tikzpicture}
\caption{Graph of the $n$-th term against $n$ for $x_n=(-1)^n+1/n$. The even points
cling near $1$, the odd points near $-1$, and the distance between the two groups is
fixed at $2$. No horizontal band contains all the terms, so this sequence is
divergent.}
\label{fig:divergen-ayun}
\end{figure}

In everyday language: the highest value still visited by the sequence as often as one
likes really is visited by some of its terms, not merely approached from afar.

\begin{teorema}
For a bounded sequence $(x_n)$, its limit superior is the largest subsequential limit:
$\limsup x_n$ equals the supremum of the set of all limits of convergent subsequences,
and that value itself is attained by some subsequence.
\end{teorema}

\begin{proof}
Write $S=\limsup x_n$. First, every subsequential limit is at most $S$. If
$x_{n_k} \to L$, then for every $k$ we have $x_{n_k} \le s_{n_k}$, and because
$n_k \ge k$ and $(s_n)$ is decreasing, $s_{n_k} \le s_k$. So $x_{n_k} \le s_k$, and
taking limits gives $L \le S$.

Second, we build a subsequence tending to $S$, starting from $n_0=0$. For every $k$,
because $s_n \to S$ and $(s_n)$ is decreasing, there is $M \ge k$ with $s_M<S+1/k$.
The tail $\{x_n : n \ge M\}$ has supremum $s_M$ and contains indices as large as we
like, so there is $n_k>\max(k,n_{k-1})$ with $n_k \ge M$ and $x_{n_k}>s_M-1/k$. For
that choice
\[
S - \frac1k \le s_M - \frac1k < x_{n_k} \le s_M < S + \frac1k ,
\]
so $|x_{n_k}-S|<1/k$. Hence $x_{n_k} \to S$.
\end{proof}

The consequence: equality of limit inferior and limit superior becomes the cheapest way
to check convergence without guessing the limit in advance.

\begin{akibat}
A bounded sequence $(x_n)$ converges to $L$ if and only if
$\liminf x_n = \limsup x_n = L$.
\end{akibat}

\begin{proof}
If $x_n \to L$, every subsequence also tends to $L$, so the set of subsequential limits
contains only $L$ and its supremum is $L$; in the same way its infimum is $L$.
Conversely, if both values equal $L$, then $t_n \le x_n \le s_n$ and both fences tend to
$L$, so the squeeze theorem gives $x_n \to L$.
\end{proof}

For unbounded sequences the convention $\limsup x_n=+\infty$ is used when $(x_n)$ is not
bounded above; the arithmetic laws for $\infty$ do not hold automatically, and we do not
use them here.

\section{Worked examples}
\tingkat{intermediate}

We return to the two opening sequences, this time with numbers that can be rechecked and
with the intermediate steps written down.

\begin{resep}
\begin{enumerate}
  \item Write the difference. Form $|x_n-L|$, then simplify until a single fraction in
        $n$ is left; for the example below, $|x_n-3|=5/(n+2)$.
  \item Solve for $n$. Solve $5/(n+2)<\varepsilon$ for $n>5/\varepsilon-2$, and name
        that threshold $g(\varepsilon)$.
  \item Choose $N$. Take $N=\lfloor g(\varepsilon)\rfloor$, the smallest admissible
        integer, then check the error at $n=N$ and $n=N+1$ before writing anything.
  \item Close with the sentence of the definition: for every $\varepsilon>0$ that $N$
        is chosen for, every $n>N$ then satisfies $|x_n-3|<\varepsilon$.
\end{enumerate}
Order warning: $\varepsilon$ is named first, $N$ answers. Setting up $N$ and then
looking for a matching $\varepsilon$ reverses the game; the result is usually an $N$
that depends on $n$, not on $\varepsilon$, and at that point the proof collapses.
\end{resep}

For $x_n=(3n+1)/(n+2)$,
\[
|x_n-3| = \left|\frac{3n+1-3n-6}{n+2}\right| = \frac{5}{n+2}.
\]
The inequality $5/(n+2)<\varepsilon$ is equivalent to $n>5/\varepsilon-2$. Because the
condition is strict, the largest integer still below that threshold is already
admissible: $N=\lfloor 5/\varepsilon-2 \rfloor$ is the smallest choice that satisfies the
definition. For $\varepsilon=10^{-2}$ the threshold is $498$ and $N=498$. Test the
boundary: at $n=498$ the error is $5/500=0.01$, which is \emph{not} smaller than
$\varepsilon$; at $n=499$ the error is $5/501=0.00998004$, which just meets it. For
$\varepsilon=10^{-4}$ we obtain $N=49998$, with error $0.0001$ at $n=49998$ and
$9.9998 \times 10^{-5}$ at $n=49999$.

The terms move very slowly: $n=10$ gives $2.583333$, $n=1000$ gives $2.995010$, and an
error of $10^{-4}$ is only reached at term $49999$. Four digits before the decimal point
are already correct long before the error is small.

\kertas{For $x_n=(3n+1)/(n+2)$, compute a valid $N$ for $\varepsilon=10^{-3}$ with the
formula $N=\lfloor 5/\varepsilon-2 \rfloor$, then check the error at $n=N$ and $n=N+1$.
Explain why the error at $n=N$ may be exactly equal to $\varepsilon$ while the definition
is still satisfied.}

For $y_n=(1+1/n)^n$, 40-digit arithmetic gives error $0.124539$ at $n=10$, $0.00135790$
at $n=10^3$, and $1.35914 \times 10^{-7}$ at $n=10^7$. The product of $n$ with its error
is nearly constant: $1.245394$, $1.357896$, then creeping to $1.3591409 = e/2$. So the
error falls like $1/n$, not $1/n^2$; one extra digit after the decimal point demands an
$n$ ten times larger. This is where the machine table at the opening of the chapter
breaks: at $n=10^7$ float64 reports $1.34327 \times 10^{-7}$, because a small relative
error in $1+1/n$ is amplified by the power $n$.

For the recursion $x_{n+1}=\sqrt{2+x_n}$ with $x_1=1$, the error relative to $2$ is
$2.679 \times 10^{-1}$ at $n=2$, $1.711 \times 10^{-2}$ at $n=4$, and
$2.615 \times 10^{-7}$ at $n=12$. The ratio of two consecutive errors settles at
$0.250000$.

The number $1/4$ is no coincidence. From the relation,
\[
x_{n+1}-2 = \frac{(2+x_n)-4}{\sqrt{2+x_n}+2}
= \frac{x_n-2}{\sqrt{2+x_n}+2},
\]
so
\[
|x_{n+1}-2| = \frac{|x_n-2|}{\sqrt{2+x_n}+2}
\le \frac{|x_n-2|}{2+\sqrt{2}} = 0.292893\,|x_n-2| ,
\]
because $x_n \ge 1$ gives $\sqrt{2+x_n} \ge \sqrt{3}$. The bound $0.292893$ holds for
all $n$, and the actual ratio tends to $1/4$ because the denominator tends to $4$. The
practical consequence: error $10^{-6}$ is reached at $n=12$, $10^{-9}$ at $n=17$, and
machine precision around $n=27$; far cheaper than bisecting an interval again and again.

\section{Computational lab: four checks}
\tingkat{advanced}

\kamus{truncation error}{galat pemotongan}{the difference between a portrait and the
actual face}

Four claims in this chapter can be checked with one short script: the explicit threshold
$N$, the error rate of $(1+1/n)^n$, machine precision, and the error ratio of the
recursion.

\begin{lab}
\begin{verbatim}
import math, mpmath as mp

def N_for(eps):
    N = int(5/eps - 2)
    while 5/(N+2) >= eps: N += 1
    while N > 0 and 5/(N+2) < eps: N -= 1
    return N

print("== 1. explicit N for (3n+1)/(n+2) -> 3 ==")
for eps in (1e-2, 1e-4):
    N = N_for(eps)
    print("eps=%.0e  N=%d  err(N)=%.6g  err(N+1)=%.6g"
          % (eps, N, abs((3*N+1)/(N+2)-3), abs((3*N+4)/(N+3)-3)))

print("== 2. (1+1/n)^n in float64 ==")
for k in range(1, 8):
    n = 10**k
    print("n=1e%d  %.16f  e-y=%.3e" % (k, (1+1/n)**n, math.e-(1+1/n)**n))

print("== 3. true error (mpmath, 40 digits) ==")
mp.mp.dps = 40
for k in (5, 6, 7):
    n = 10**k
    y = (1+mp.mpf(1)/n)**n
    print("n=1e%d  e-y=%s  (float64: %.3e)"
          % (k, mp.nstr(mp.e-y, 6), math.e-(1+1/n)**n))

print("== 4. recursion x_{n+1}=sqrt(2+x_n), x_1=1 ==")
x, xs = 1.0, []
for _ in range(30):
    xs.append(x); x = math.sqrt(2+x)
for n in (1, 2, 3, 4, 5, 6, 11, 12, 13, 14):
    e = abs(xs[n-1]-2)
    r = "" if n == 1 else "  ratio=%.6f" % (e/abs(xs[n-2]-2))
    print("n=%2d  x=%.12f  error=%.3e%s" % (n, xs[n-1], e, r))
e6 = next(n for n in range(1, 31) if abs(xs[n-1]-2) < 1e-6)
e9 = next(n for n in range(1, 31) if abs(xs[n-1]-2) < 1e-9)
print("error < 1e-6 first at n=%d ; error < 1e-9 at n=%d" % (e6, e9))
\end{verbatim}
\end{lab}

\newpage
The output, as it stands:

\begin{verbatim}
== 1. explicit N for (3n+1)/(n+2) -> 3 ==
eps=1e-02  N=498  err(N)=0.01  err(N+1)=0.00998004
eps=1e-04  N=49998  err(N)=0.0001  err(N+1)=9.9998e-05
== 2. (1+1/n)^n in float64 ==
n=1e1  2.5937424601000023  e-y=1.245e-01
n=1e2  2.7048138294215285  e-y=1.347e-02
n=1e3  2.7169239322355936  e-y=1.358e-03
n=1e4  2.7181459268249255  e-y=1.359e-04
n=1e5  2.7182682371922975  e-y=1.359e-05
n=1e6  2.7182804690957534  e-y=1.359e-06
n=1e7  2.7182816941320818  e-y=1.343e-07
== 3. true error (mpmath, 40 digits) ==
n=1e5  e-y=1.35913e-5  (float64: 1.359e-05)
n=1e6  e-y=1.35914e-6  (float64: 1.359e-06)
n=1e7  e-y=1.35914e-7  (float64: 1.343e-07)
== 4. recursion x_{n+1}=sqrt(2+x_n), x_1=1 ==
n= 1  x=1.000000000000  error=1.000e+00
n= 2  x=1.732050807569  error=2.679e-01  ratio=0.267949
n= 3  x=1.931851652578  error=6.815e-02  ratio=0.254333
n= 4  x=1.982889722748  error=1.711e-02  ratio=0.251074
n= 5  x=1.995717846477  error=4.282e-03  ratio=0.250268
n= 6  x=1.998929174953  error=1.071e-03  ratio=0.250067
n=11  x=1.999998954179  error=1.046e-06  ratio=0.250000
n=12  x=1.999999738545  error=2.615e-07  ratio=0.250000
n=13  x=1.999999934636  error=6.536e-08  ratio=0.250000
n=14  x=1.999999983659  error=1.634e-08  ratio=0.250000
error < 1e-6 first at n=12 ; error < 1e-9 at n=17
\end{verbatim}

Three things need to be read from that output. First, the error at $n=N$ is exactly
equal to $\varepsilon$ in both lines; that is why the definition uses the strict
condition $n>N$. Second, the $e-y$ column shrinks by a factor of ten each time $n$ grows
tenfold up to $n=10^6$, then stops following the pattern; part 3 shows that what stops
is float64 precision, not the convergence of the sequence. The truncation error keeps
shrinking like $e/(2n)$, but the rounding error takes over as soon as the two meet.
Third, the error ratio in part 4 falls to $0.250000$ and stays there, exactly as in the
hand computation before. The program proves nothing: it only checks claims we have
already proved.

\begin{untukdosen}
Teaching order: start with the $\varepsilon$-$N$ game in class, not with the theorems.
Point to two rows of students; one names $\varepsilon$, one answers $N$ for
$x_n=(3n+1)/(n+2)$, and the instructor writes it down until the formula appears by
itself. After that come uniqueness, the algebraic properties, the squeeze theorem,
monotone sequences, then subsequences and Cauchy. The limit superior is placed at the
end, as ``a limit for sequences that have no limit''.

Two typical misconceptions. First, students think $N$ may depend on $n$, then write
``for every $n$ there is $N=n$'' and believe the definition is met; the symptom shows
when they cannot produce a single $N$ for the whole tail. Second, they identify ``terms
getting closer together'' with ``convergent'' and use the Cauchy criterion without
checking the space; the decimal truncations of the square root of two inside
$\mathbb{Q}$ are an effective cure.

Board activity: ask students to write the negation of the definition of a limit, then
prove the divergence of $(-1)^n$ with that negation, not by contradiction. The class is
ready to move on when they can produce an explicit $N$ without being given a formula
and name one bounded sequence that is divergent.
\end{untukdosen}

\section{Summary}
\tingkat{basic}

\begin{ringkasan}
\begin{itemize}
  \item $x_n \to L$ means that for every $\varepsilon>0$ there is $N$ such that
        $|x_n-L|<\varepsilon$ for all $n>N$. The $\varepsilon$ is chosen first; $N$
        answers, and may depend on $\varepsilon$ alone.
  \item The limit is unique, and a convergent sequence is bounded. The converse fails:
        boundedness is only a necessary condition.
  \item Limits respect $+$, $\times$, and division by a nonzero denominator; order is
        never jumped over, only weakly preserved.
  \item The squeeze theorem turns trapping into a limit, and the lemma $q^n \to 0$
        supplies a comparison used repeatedly. A bounded monotone sequence always
        converges, and its limit is the supremum or the infimum.
  \item Every bounded sequence has a convergent subsequence; a real sequence converges
        if and only if it is Cauchy, and that equivalence is completeness.
  \item $\liminf \le \limsup$, and the two are equal to $L$ exactly when the sequence
        converges to $L$.
\end{itemize}
\end{ringkasan}

\section{Exercises}
\tingkat{intermediate}

\begin{latihan}
\begin{enumerate}
  \item Prove from the definition that $x_n=(2n-1)/(3n+4)$ converges to $2/3$, and give
        an explicit $N$ as a function of $\varepsilon$. Compare it with the smallest
        $n$ that satisfies the condition.
  \item Prove: if $x_n \to L$ then $|x_n| \to |L|$. Show by example that the converse is
        false.
  \item Compute $\limsup$ and $\liminf$ of $x_n=1+(-1)^n+1/n$ and of
        $y_n=(-1)^n(1+1/n)$. Check your answers with Python for $n$ up to $1000$.
  \item Evaluate this claim: ``Because $(1+1/n)^n$ is increasing and bounded above by
        $3$, its limit is $3$.'' True, false, or not yet enough evidence? Say which part
        is right and which part is wrong.
  \item Prove that $(2^n+3^n)^{1/n} \to 3$ by writing out both fences completely,
        including a proof that $2^{1/n} \to 1$.
  \item For the recursion $x_{n+1}=\sqrt{2+x_n}$ with $x_1=1$, use Python to find the
        smallest term with error below $10^{-6}$ and below $10^{-12}$. Compare with the
        guess from the ratio $1/4$.
  \item Prove that a Cauchy sequence is bounded, without using Bolzano--Weierstrass.
        Then prove that a bounded monotone sequence is Cauchy.
  \item For $x_n=(3n+1)/(n+2)$ and $\varepsilon=10^{-3}$, determine the smallest
        admissible $N$, then compute the error at $n=N$, $N+1$, and $N+2$ with Python.
        Explain why $N-1$ is not admissible.
  \item Prove: if every subsequence of $(x_n)$ has a subsequence converging to $L$, then
        $x_n \to L$.
\end{enumerate}

Hints for the three hardest items. Item 5: use
$1 \le (1+(2/3)^n)^{1/n} \le 2^{1/n}$ and then the lemma $2^{1/n} \to 1$ via Bernoulli's
inequality. Item 7: for the second part, compare $|x_n-x_m|$ with $x_m-x_n$ when $m>n$
and use monotone convergence. Item 9: suppose $x_n$ does not tend to $L$, build a
subsequence staying a fixed distance from $L$, then use Bolzano--Weierstrass.
\end{latihan}

\begin{untukpeneliti}
What is really proved here is only the following equivalence, and all of it rests on one
axiom: the supremum property. In an ordered field, the following statements are
equivalent: every nonempty set bounded above has a supremum; every bounded increasing
sequence converges; every Cauchy sequence converges; every bounded sequence has a
convergent subsequence. The chain of proofs is in Rudin, \emph{Principles of Mathematical
Analysis}, Chapter 3 (the sections Convergent Sequences, Subsequences, Cauchy Sequences)
and Abbott, \emph{Understanding Analysis}, Chapter 2 (sections 2.2 to 2.6). What is
popular but not proved: ``a Cauchy sequence always converges''. That is false without
completeness, and the decimal truncations of the square root of two inside
$\mathbb{Q}$ are the counterexample.

An open question: how fast a sequence converges, and whether that rate can always be
computed. The answer is no. A classical result of Specker (1949) gives a computable,
increasing, bounded sequence of rationals whose limit is not computable; so
``convergent'' does not automatically give an error bound that can be programmed.

Limits of validity: this whole chapter works in $\mathbb{R}$ with the usual distance. In
a general metric space, Bolzano--Weierstrass becomes a statement about compactness, and
the Cauchy criterion remains a necessary condition but no longer sufficient. The squeeze
theorem is also silent when the two fences do not tend to the same value; there we need a
finer comparison, and that is the material of Chapter 5.
\end{untukpeneliti}

\peta{Chapter 4 in Rudin, \emph{Principles of Mathematical Analysis}: Chapter 3 (Numerical
Sequences and Series, with the sections Convergent Sequences, Subsequences, Cauchy
Sequences, Upper and Lower Limits, and Some Special Sequences). In Abbott,
\emph{Understanding Analysis}: Chapter 2 (The Limit of a Sequence, The Algebraic and
Order Limit Theorems, The Monotone Convergence Theorem and a First Look at Infinite
Series, Subsequences and the Bolzano--Weierstrass Theorem, The Cauchy Criterion).
The axioms and completeness consequences used here are built in Chapter 3 of this book,
not repeated.}
